Lo scopo di questa tesi magistrale è di derivare le equazioni effettive per un modello microscopico per la chemotassi di $N$ amebe attratte dalle molecole di cAMP che esse stesse secernono in impulsi ditribuiti nel tempo con una legge di Poisson.
L'interazione è espressa mediante un potenziale di tipo Gaussiano.
Il modello microscopico è ispirato a \cite{ana}, che, insieme a \cite{peter-boers}, è la fonte principale delle tecniche adottate nella dimostrazione del risultato principale.
Inoltre discutiamo delle somiglianze e delle differenze che la descrizione macroscopica che abbiamo derivato ha con il modello macroscopico classico per la chemotassi, cioè \emph{Keller-Segel}.